\chapter{Bernoulli Numbers}

\section{Faulhaber's Formula}

\begin{theorem}
  \begin{equation*}
    \sum_{k=1}^n k^m = \frac{1}{m+1}\sum_{r=0}^n\binom{m+1}{r} B_r n^{m+1-r}.
  \end{equation*}
\end{theorem}

The following theorem is a consequence of Faulhaber's formula, but can also be derived independently.

\begin{theorem}
  Assume that $m$ is even. Then
  \begin{equation*}
    \sum_{j=1}^{p-1} j^m\equiv -\frac{B_{m+1}}{m+1} p\pmod{p^2}.
  \end{equation*}
  Equivalently,
  \begin{equation*}
    \frac{1}{p}\sum_{j=1}^{p-1} j^m\equiv -\frac{B_{m+1}}{m+1}\pmod{p}.
  \end{equation*}
\end{theorem}

\section{Fundamental Identities Involving Bernoulli Numbers}

\begin{equation*}
  \frac{t}{e^t-1}=\sum_{n=0}^{\infty}\frac{B_n}{n!} t^n.
\end{equation*}

\begin{equation*}
  \sum_{k=0}^n\binom{n}{k}B_k = B_n+nB_{n-1}
\end{equation*}

\begin{equation*}
  \sum_{k=0}^n\binom{n}{k}2^kB_k = (2-2^n)B_n  
\end{equation*}

\begin{equation*}
  \sum_{k=1}^n\binom{n}{k} B_k B_{n-k}=-n(B_n - B_{n-1})
\end{equation*}
https://math.stackexchange.com/questions/3243169/euler-sum-with-bernoulli-numbers

\begin{equation*}
  \sum_{k=j}^n\frac{1}{k}\binom{n}{k-1}\binom{k}{j} B_{k-j}=\delta_{nj}
\end{equation*}
where $\delta_a$ is the Kronecker $\delta$-function.

\section{Values of the \texorpdfstring{$\zeta$}{ζ}-function}

\begin{equation*}[Euler]
  \zeta(2n)=\sum_{k=1}^{\infty}\frac{1}{k^{2n}}=(-1)^{n+1}\frac{(2\pi)^{2n}}{2(2n)!}B_{2n}.
\end{equation*}

\section{Kummer's Congruence for Bernoulli Numbers}

If $n\equiv m\pmod{p-1}$ and $m$ and $n$ are positive even integers, then
\begin{equation*}
  \frac{B_n}{n}\equiv \frac{B_m}{m}\pmod{p}.
\end{equation*}
This has applications to $p$-adic L-functions. More generally, if $n\equiv m\pmod{\phi(p^k)}$ and $m$ and $n$ are positive even integers, then
\begin{equation*}
  (1-p^{n-1})\frac{B_n}{n}\equiv (1-p^{m-1})\frac{B_m}{m}\pmod{p^k}.
\end{equation*}

This leads to some observations about the structure of cyclotomic fields and class groups, and eventually to Iwasawa theory.

\section{Congruences Related to Bernoulli Numbers}

\begin{definition}
  A prime $p$ is said to be \emph{regular} if it does not divide the numerator of any of the Bernoulli numbers $B_{2k}$ for $2k\leq p-3$.
\end{definition}

\begin{theorem}[The von~Staudt--Clausen Theorem]
  Let  $B_k$ denote the $k$\supth{} Bernoulli number. Then
  
  \begin{equation*}
    B_{2n}+\sum_{p-1\mid 2n}\frac{1}{p}
  \end{equation*}
  is an integer.
\end{theorem}

\begin{theorem}[Glaisher's Congruence]
  For $p>3$,
  \begin{equation*}
    \binom{2p-1}{p-1}\equiv 1-\frac{2}{3}p^2 \mathrm{num}(B_{p-3})\pmod{p^3}
  \end{equation*}
\end{theorem}

https://scispace.com/pdf/notes-on-glaisher-s-congruences-25n3c5gnc3.pdf \\
https://arxiv.org/pdf/1301.4888

\section{Fermat Quotients}

Let $q_p(a):= \frac{a^{p-1}-1}{p}$ be the \emph{Fermat quotient}.

\begin{theorem}[Vandiver's Criterion]
  If $p\mid q_p(a)$, then $p$ divides the numerator of the $(p-1)$\supth{} Bernoulli number $B_{p-1}$.
\end{theorem}

Some further interesting results about Fermat quotients:
\begin{equation*}
  q_p(ab)\equiv q_p(a)+q_p(b)\pmod{p}.
\end{equation*}
Leads to Wieferich primes and cyclotomic class groups.

\begin{theorem}[Skula's Congruence]
  \begin{equation*}
    q_p(2)^2\equiv \sum_{k=1}^{p-1}\frac{1}{k^2}\pmod{p}.
  \end{equation*}
\end{theorem}

There are some generalizations to Euler quotients too:
\begin{equation*}
  q_n(a):=\frac{a^{\phi(n)}-1}{n}
\end{equation*}
for $a$ relatively prime to $n$.

%%% Local Variables:
%%% mode: latex
%%% TeX-master: "elementary-number-theory"
%%% TeX-engine: luatex
%%% End:
